\documentclass[10pt,a4paper]{amsart}
\usepackage[margin=19mm]{geometry}
\usepackage{amsmath,amssymb,mathtools}
\usepackage[scr=rsfs,cal=euler]{mathalfa}
\usepackage{xfrac}
\usepackage[unicode,hidelinks,bookmarks=false]{hyperref}
\newcommand{\ie}{i.e.\ }
\newcommand{\cf}{cf.\ }
\newcommand{\resp}{respectively\ }
\newcommand{\Subsection}{\subsection}
\providecommand{\nobreakdash}{\nobreak\hbox{-}}
\setlength{\emergencystretch}{2em}
\multlinegap0pt
\usepackage{bm}
\usepackage{bbm}
\usepackage{dsfont}
\usepackage{xspace}
\usepackage{cancel}
\usepackage{mathtools}
\usepackage{amsthm}
\makeatletter
\@ifundefined{theorem}{\newtheorem{theorem}{Theorem}}{}
\makeatother
\title[On a Generalized Compartment Model for Ethanol Metabolism in]{On a Generalized Compartment Model for Ethanol Metabolism in the Human Body}
\author{Manh Tuan Hoang, Thi Kim Quy Ngo, Benjamin Wacker}
\date{Source: arXiv:2606.30439v2 (CC BY 4.0)}
\begin{document}
\maketitle

\noindent\emph{Source and licence.} On a Generalized Compartment Model for Ethanol Metabolism in the Human Body. Version arXiv:2606.30439v2 (CC BY 4.0), distributed under the Creative Commons Attribution 4.0 International licence, as recorded in the arXiv licence metadata for this exact version and shown on the abstract page https://arxiv.org/abs/2606.30439v2. Full rights notice: licenses/arxiv-2606.30439-CC-BY-4.0.txt.\par\smallskip

\noindent\textit{Editorial scope.} Counted unit: the Theorem whose source label is \texttt{Theorem5new} in the article (source0.tex lines 552--554, source label \texttt{Theorem5new}), delivered together with the article's own complete original proof of it (source0.tex lines 556--635, one \texttt{proof} environment of 80 lines). No mathematics is written, repaired or re-derived: statement and proof are the article's own text line for line. Required context retained uncounted: eq:DE1 (lines 277--288); eq:DE2 (lines 424--426); eq:DE7 (lines 432--443); Theorem5 (lines 488--490); eq:DE3a (lines 546--548). Every definition, display and label of those lines is retained unchanged. Omitted from this package and not redistributed: 1--276, 289--423, 427--431, 444--487, 491--545, 549--551, 555--555, 636--1019. Nothing inside a retained excerpt is omitted or altered: every retained body is delivered exactly as the article's lines are written, so no omission receipt of any kind accompanies this package. Citations: the retained excerpts carry no citation command at all, so no bibliography entry of the article is transcribed and no reference is redistributed. Reference adaptation: none was needed and none was made; every reference inside the delivered unit resolves inside it, so no unresolved reference or citation is produced. Printed slips kept literal: no printed word, symbol, hypothesis or number of the counted statement or of its proof was altered. Layout and class: the article's own class is \texttt{article} with packages including fontenc, authblk, amsmath, amsfonts, amssymb, ulem, lineno, xcolor, amsthm, graphicx, marginnote, subcaption, hyperref, geometry; the wrapper is the lane's pinned \texttt{amsart} editorial wrapper at 10pt on A4, which restates the theorem-like environments the retained text needs and adds the article's own macro definitions that the retained text needs (none), with extra wrapper packages (bm, bbm, dsfont, xspace, cancel, mathtools). The delivered numbering is the unit's own. Assets: no retained span contains a figure environment, a graphics-inclusion command, a PDF-page inclusion, a table or a TikZ picture; no image file is copied into this package, the package declares no asset dependency and no image file is redistributed with it. Licence: version arXiv:2606.30439v2 of this article is distributed under the Creative Commons Attribution 4.0 International licence, as recorded in the arXiv licence metadata for this exact version and shown on the abstract page https://arxiv.org/abs/2606.30439v2. A full rights notice is delivered at licenses/arxiv-2606.30439-CC-BY-4.0.txt. Characters: every retained excerpt is pure ASCII, so the delivered engine, XeLaTeX with its default Latin Modern text font, reports no missing character.
\begin{equation}\label{eq:DE1}
\begin{split}
A_{n + 1} &= A_n  -a\Delta t A_n, \quad A_0 > 0,\\
%%
%%
%%
B_{n+1} &= B_n + b\Delta t C_n - b\Delta t B_n + a\Delta t A_n,\quad B_0 \geq 0,\\
%%
%%
C_{n + 1} &= C_n +  b\Delta t B_n- b\Delta t C_n- \Delta tf(C_n), \quad C_0 \geq 0,
\end{split}
\end{equation}

\begin{equation}\label{eq:DE2}
\Delta t \leq \min\bigg\{\dfrac{1}{a},\,\,\,\dfrac{1}{b + M}\bigg\}.
\end{equation}

\begin{equation}\label{eq:DE7}
\begin{split}
A_{n + 1} &= (1 -a\Delta t) A_n,\\
%%
%%
%%
B_{n+1} &= (1 - b\Delta t)B_n + b\Delta t C_n + a\Delta t A_n,\\
%%
%%
C_{n + 1} &= \big(1 - b\Delta t - g(C_n)\Delta t\big)C_n +  b\Delta t B_n.
\end{split}
\end{equation}

\begin{theorem}[Local asymptotic stability analysis]\label{Theorem5}
Under the condition \eqref{eq:DE2}, the trivial equilibrium point of \eqref{eq:DE1} is locally asymptotically stable for all parameter values.
\end{theorem}

\begin{equation}\label{eq:DE3a}\
\min_{0 \leq C \leq S_0}g(C) = m > 0,
\end{equation}

\begin{theorem}[Global asymptotic stability analysis]\label{Theorem5new}
The trivial equilibrium point of \eqref{eq:DE1} is not only locally asymptotically stable but also globally asymptotically stable for all parameter values if \eqref{eq:DE2} and \eqref{eq:DE3a} hold.
\end{theorem}

\begin{proof}
From Theorem \ref{Theorem5}, it is sufficient to show that
\begin{equation*}
\lim_{n \to \infty}\big(A_n,\,B_n,\,C_n\big) = (0,\,0,\,0).
\end{equation*}
%%
It follows from the last equation of \eqref{eq:DE7} that
\begin{equation*}
C_{n + 1} = (1 - b\Delta t - g(C_n)\Delta t)C_n +  b\Delta t B_n \leq (1 - b\Delta t - m\Delta t)C_n +  b\Delta t B_n.
\end{equation*}
Thus, \eqref{eq:DE7} implies that
\begin{equation*}
\begin{pmatrix}
A_{n + 1}\\
B_{n + 1}\\
C_{n + 1}
\end{pmatrix}
\leq
P
\begin{pmatrix}
A_n\\
B_n\\
C_n
\end{pmatrix},
\end{equation*}
where $P$ is defined by
\begin{equation*}
P :=
\begin{pmatrix}
1 - a\Delta t&0&0\\
%%
%%
a\Delta t&1 - b\Delta t& b\Delta t\\
%%
%%
0&b\Delta t & 1 - b\Delta t - m\Delta t.
\end{pmatrix}.
\end{equation*}
From \eqref{eq:DE2}, all the entries of $P$ are positive. Thus, for $n \geq 1$
\begin{equation}\label{eq:DE8}
\begin{pmatrix}
A_n\\
B_n\\
C_n
\end{pmatrix}
\leq P^n
\begin{pmatrix}
A_0\\
B_0\\
C_0
\end{pmatrix}.
\end{equation}
The matrix $P$ always has an eigenvalue given by $\mu_1 = 1 - a\Delta t$, which satisfies $|\mu_1| < 1$. Meanwhile, the two remaining eigenvalues coincide with those of the submatrix
%
\begin{equation*}
P^* :=
\begin{pmatrix}
1 - b\Delta t& b\Delta t\\
%%
%%
b\Delta t & 1 - b\Delta t - m\Delta t.
\end{pmatrix}.
\end{equation*}
%
Under the condition \eqref{eq:DE2}, we have the estimates
%
\begin{equation*}
\begin{split}
\det(P^*) - 1 &= -b\Delta t - m\Delta t - b\Delta t(1 - m\Delta t) < 0\\
%%
%%
1 - Tr(P^*) + \det(P^*) &=  bm(\Delta t)^2> 0,\\
%%
%%
1 + Tr(P^*) + \det(P^*) &= 1 - 4b\Delta t - 2m\Delta t +  bm(\Delta t)^2\\
&> 4\big(1 - b\Delta t - m\Delta t\big) + bm(\Delta t)^2 > 0.
\end{split}
\end{equation*}
Consequently, we conclude that all the eigenvalues $\mu_i$ ($i = 1, 2, 3$) of $P$ satisfy $|\mu_i| < 1$, which implies that $\lim_{n \to \infty}P^n = 0$. Then, it follows from \eqref{eq:DE8} that $\lim_{n \to \infty}\big(A_n,\,B_n,\,C_n\big) = (0,\,0,\,0)$. This is the desired conclusion. The proof is complete.
\end{proof}

\end{document}
